\chapter{Coccidioidomycosis (Valley Fever)}
\index{Coccidioidomycosis (Valley Fever)}
\label{sec:Coccidioidomycosis}

\section{Overview}
Coccidioidomycosis (Valley fever) is a systemic fungal infection caused by \textit{Coccidioides immitis} and \textit{C. posadasii}, dimorphic fungi endemic to the Sonoran Desert (southwestern US, northern Mexico, and parts of Central and South America), with 150,000 infections annually in the US, and incidence increasing due to climate change.
\section{Causes}
This condition happens because of inhalation of arthroconidia leading to conversion into spherules containing endospores in the lung, with a robust Th1 response controlling infection.     
\section{Risk Factors}

\begin{itemize}[noitemsep]
  \item Genetic predisposition and family history
  \item Advancing age
  \item Lifestyle factors (diet, exercise, smoking, alcohol)
  \item Environmental exposures
  \item Underlying medical conditions
  \item Medication use
\end{itemize}
\section{Signs and Symptoms}

\subsection{Common symptoms}
\begin{itemize}[noitemsep]
  \item Pain or discomfort in the affected area
  \end{itemize}

\subsection{Emergency warning signs}
\begin{itemize}[noitemsep]
  \item Severe or worsening symptoms of coccidioidomycosis (valley fever) require immediate medical evaluation
\end{itemize}
\section{Progression and Stages}
The condition may progress through different stages of severity over time. Early stages often involve mild or subtle symptoms that may be overlooked. Moderate stages involve more noticeable symptoms affecting daily function. Advanced stages may involve significant impairment and complications.
\section{Complications}
You may experience acute lung: 60\% asymptomatic, 40\% develop Valley fever (influenza-like illness with fever, cough, chest pain, joint pain, and redness of the skin nodosum)

\begin{itemize}[noitemsep]
  \item Chronic lung: persistent cough, hemoptysis, cavities
  \item Disseminated coccidioidomycosis: skin, bones/joints, and CNS (coccidioidal meningitis---the most feared complication).
  \item Disease progression leading to worsening symptoms
  \item Secondary infections or injuries
  \item Side effects of treatment
  \item Reduced quality of life
  \item Emotional and psychological impact
\end{itemize}
\section{Diagnosis}
Doctors diagnose this using culture, serology, and histopathology. Comprehensive medical history and physical examination Laboratory tests to assess organ function and rule out other conditions Imaging studies to visualize affected structures Specialized testing depending on the affected system Consultation with relevant specialists Monitoring of disease progression over time

\begin{itemize}[noitemsep]
  \item Comprehensive medical history and physical examination
  \item Laboratory tests to assess organ function
  \item Imaging studies to visualize affected structures
  \item Specialized testing depending on the affected system
  \item Consultation with relevant specialists
\end{itemize}
\section{Prevention}
\begin{itemize}[noitemsep]
  \item Maintain a healthy lifestyle with balanced diet and exercise
  \item Avoid known risk factors
  \item Attend regular medical check-ups for early detection
  \item Follow recommended screening guidelines
\end{itemize}
\section{Treatment Options}
Treatment options include fluconazole for mild-moderate disease and liposomal amphotericin B for severe disseminated disease. Medications to manage symptoms and address underlying mechanisms Lifestyle modifications and self-care measures Physical therapy and rehabilitation Surgical intervention when indicated Regular monitoring and follow-up care Patient education and support services

\begin{itemize}[noitemsep]
  \item Medications to manage symptoms and address underlying mechanisms
  \item Lifestyle modifications and self-care measures
  \item Physical therapy and rehabilitation
  \item Surgical intervention when indicated
  \item Regular monitoring and follow-up care
\end{itemize}
\section{Living with It}
\begin{itemize}[noitemsep]
  \item Follow your treatment plan as prescribed
  \item Maintain regular follow-up with your healthcare provider
  \item Make necessary lifestyle adjustments
  \item Seek support from family, friends, or support groups
  \item Stay informed about your condition
\end{itemize}
\section{Outlook}
\begin{itemize}[noitemsep]
  \item Most people recover well for acute lung disease
  \item Meningitis requires lifelong therapy.
\end{itemize}
